\chapter{\texorpdfstring{$\Mgn$}{Stable pointed curves} is a stack}\label{ch:curves}

\section{A family of stable pointed curves}

Start with a scheme $S$ over $\CC$, and fix integers $g,n\geq0$ with
\begin{equation}\label{eq:stable-range}
                         2g-2+n>0.
\end{equation}
We want to put one complete, connected, possibly nodal curve over each
point of $S$, together with $n$ labelled points that vary as sections.
The scheme $S$ is the parameter scheme; it is not the curve being
classified. The inequality is the \emph{stable range}: for example,
$(g,n)=(2,0),(1,1),(0,3)$ are allowed, and $(0,0)$ is not allowed
for stable pointed curves.

\begin{Def}[Family of stable pointed curves]\label{def:stable-family}
The input is a $\CC$-scheme $S$. A family over $S$ consists of a scheme
$C$, a proper flat finitely presented morphism $\pi:C\to S$, and
sections $\sigma_a:S\to C$ for $1\leq a\leq n$, such that:
\begin{enumerate}
\item For every geometric point $\bar s:\Spec\Omega\to S$, where
$\Omega$ is algebraically closed, the fiber
$C_{\bar s}=C\times_S\Spec\Omega$ is a connected reduced curve with
only ordinary nodes as possible singularities, and arithmetic genus
$g$.
\item Each section takes values in the smooth locus of $\pi$, and the
images of different sections are disjoint.
\item In the normalization of each geometric fiber, every component
of genus zero carries at least three special points, and every
component of genus one carries at least one special point. A special
point is a marked point or a preimage of a node; a self-node contributes
two such preimages to the same component.
\end{enumerate}
We write $\xi=(S,C,\pi,\sigma_1,\ldots,\sigma_n)$ for all these data.
\end{Def}

Here ``finitely presented'' rules out infinitely many equations or
parameters in the relative construction; it is the useful finiteness
condition when $S$ need not be noetherian. Flatness is a condition on
the family, not just on its individual fibers. Properness is the
algebraic completeness condition. The sections satisfy
$\pi\sigma_a=\id_S$ and remember their labels. Their disjointness is
a condition on the entire family, including nonreduced bases.

At a node of a geometric fiber, the completed local ring is
$\Omega[[x,y]]/(xy)$: two smooth branches meet. The arithmetic genus
of a proper connected reduced curve over $\Omega$ is
$g=\dim_\Omega H^1(C_{\bar s},\mathcal O_{C_{\bar s}})$.
Thus a genus-two fiber need not be a smooth surface with two handles;
it can be reducible, with the genus distributed among its components
and its incidence graph.

The normalization separates the branches at each node. If a normalized
component $E$ has genus $h$ and $r$ special points, the stability test
can be written as $2h-2+r>0$. In particular, its content concerns
\emph{each component}, not just the total genus. The geometric facts
relating this test to automorphisms and line bundles are explained in
Alper \cite[Definition 6.3.2, Proposition 6.3.5, pp.~225--226]{Alper}.
The family definition is \cite[Definition 6.3.20, p.~231]{Alper}.

\begin{Ex}[A smooth genus-two object]\label{ex:genus-two}
Let $C_0$ be the smooth projective completion of
\[
                            y^2=x^6-1
\]
over $\CC$. More precisely it is the smooth double cover of $\PP^1$
branched at the six distinct sixth roots of unity. The double-cover
genus formula gives $2g(C_0)-2=2(-2)+6=2$, hence $g(C_0)=2$.
The affine equation alone would not be a proper curve, which is why
the completion is part of the description.

With no markings, $C_0\to\Spec\CC$ is an object for $(g,n)=(2,0)$.
Its only normalized component has $h=2,r=0$, so it is stable.
The involution $\iota(x,y)=(x,-y)$ extends over the completion.
Stability allows this nontrivial automorphism; it asks for a finite
automorphism group, not a trivial group. If we want one marking,
$p=(1,0)$ is a smooth point fixed by $\iota$, so
$(C_0,p)$ is a useful example for $(2,1)$ too.
\end{Ex}

\begin{Ex}[A stable nodal genus-two object]\label{ex:genus-two-nodal}
Take smooth elliptic curves $E_1,E_2$ with chosen points $q_1,q_2$,
and identify these points to one ordinary node $q$. The result
$C=E_1\cup_q E_2$ is projective and connected. The normalization
sequence is
\[
0\longrightarrow\mathcal O_C\longrightarrow
\mathcal O_{E_1}\oplus\mathcal O_{E_2}
\longrightarrow\mathbb C_q\longrightarrow0,
\]
where the last map takes the difference of values at the two branches.
It gives $\chi(\mathcal O_C)=0+0-1=-1$, hence genus two.
Each normalized elliptic component has its one node preimage as a
special point. Thus both have $2h-2+r=1>0$. This is a stable
unmarked curve even though neither elliptic component would be
stable as an isolated unmarked curve.
\end{Ex}

\begin{Ej}\label{ex:stability-test}
A rational component has one self-node and one additional marked
smooth point. How many special points are on its normalization?
Does this component pass the stability test? Why would counting the
self-node just once give the wrong answer?
\end{Ej}

\subsection{Base and notation for the whole construction}

Our base category is $\Sch/\CC$, the category of all $\CC$-schemes
and $\CC$-morphisms, with the \emph{\'etale topology}. A cover is a
family of \'etale maps $\{S_i\to S\}$ whose images together cover
$S$. Ordinary open covers are examples, but a cover need not consist
of inclusions. For example,
\[
\Spec\CC[u,u^{-1}]\longrightarrow\Spec\CC[t,t^{-1}],\qquad t=u^2,
\]
is a one-map \'etale cover. Locally it supplies a choice of square
root; globally there is no preferred choice.

We use $\Mgn$ for the stable pointed-curve construction throughout.
The notation $\mathcal M_{g,n}$ denotes its open substack of
\emph{smooth} curves, once the stack has been constructed. With
$n=0$, write $\overline{\mathcal M}_g$ and $\mathcal M_g$,
respectively. An informal ``$M_g$'' must be disambiguated: smooth
curves and their stable compactification are different moduli
problems. Plain noncalligraphic symbols in some references instead
denote coarse spaces; we do not use them for our stacks.

\section{First construct an ordinary category}\label{sec:curve-category}

\subsection{Objects and arrows}

The objects of our \emph{total category} are all the tuples $\xi$
of Definition~\ref{def:stable-family}, allowing $S$ to vary. Thus
the base is part of the object. The same curve family pulled back
to a new base is a new object, even if its fibers look familiar.

\begin{Def}[An arrow between families]\label{def:curve-arrow}
Given objects $\xi=(C\xrightarrow{\pi}S,\sigma_\bullet)$ and
$\eta=(D\xrightarrow{\rho}T,\tau_\bullet)$, an arrow
$(u,h):\xi\to\eta$ consists of maps $u:S\to T$ and $h:C\to D$
such that
\[
\begin{tikzcd}[row sep=large,column sep=large]
C\arrow[r,"h"]\arrow[d,"\pi"']&D\arrow[d,"\rho"]\\
S\arrow[r,"u"']&T
\end{tikzcd}
\qquad
\rho h=u\pi,\qquad h\sigma_a=\tau_a u\quad(1\leq a\leq n),
\]
and the square is cartesian. Explicitly, the induced map
\[
             (\pi,h):C\longrightarrow S\times_TD
\]
must be an isomorphism.
\end{Def}

The input is two families. The data of an arrow say that the first
is the second viewed over the new base $S$, with an identification
of the pulled-back total space. The condition is an isomorphism of
\emph{schemes over $S$}, not merely a bijection of geometric fibers.
The map $h:C\to D$ itself need not be an isomorphism when $u$ is
not one. For instance, inclusion of a fiber is usually not an
isomorphism of total spaces.

Why not allow arbitrary maps between curves? The moduli problem
identifies two presentations of the same family. A degree-two map
between curves or a constant map changes the geometric object; it
is not such an identification. In particular, a constant
$C_0\to C_0$ is not an arrow here. This choice of arrows is essential
to obtaining groupoids as fibers. It does not mean that every arrow
of the total category is invertible.

\begin{Ex}[Arrows over a point]
Over $\Spec\CC$, the only base endomorphism as a $\CC$-scheme is
the identity. An arrow from $C_0$ to a second genus-two curve $D_0$
is therefore exactly a curve isomorphism $h:C_0\xrightarrow{\sim}D_0$.
If points are marked, $h$ must send each labelled point to the
point with the same label. The identity and $\iota$ give different
arrows $C_0\to C_0$, including in the one-pointed example $(C_0,(1,0))$.
An object and its automorphisms are retained, not replaced by just
one element of a set of isomorphism classes.
\end{Ex}

\subsection{Identities, composition, and the axioms}

The identity of $\xi$ is $(\id_S,\id_C)$. Its square is cartesian
and it preserves the sections. If $(u,h):\xi\to\eta$ and
$(v,k):\eta\to\zeta$, where $\zeta=(E\to U,\upsilon_\bullet)$,
define
\[
                        (v,k)\circ(u,h)=(vu,kh).
\]
We must check that this is still an allowed arrow. The base square
commutes because the two smaller squares do. The section equation is
$kh\sigma_a=k\tau_a u=\upsilon_a vu$. Finally,
\[
C\simeq S\times_TD
\simeq S\times_T(T\times_UE)
\simeq S\times_UE,
\]
so the composite square is cartesian. Composition is associative
because scheme morphisms compose associatively in both entries of
the pair. Composing with the stated identities changes neither
entry. We have now checked all the category axioms.

\MCcheckpoint{So far we have an ordinary category. There has been no
gluing theorem and no assertion about a parameter space representing it.}

\begin{Ej}\label{ex:total-vs-fiber}
Let $C_0\times\mathbb A^1\to\mathbb A^1$ be the constant
genus-two family. Explain why the inclusion of the fiber at $0$
defines an arrow in the total category, but has no inverse there.
Why does this not contradict the forthcoming claim that each fiber
category is a groupoid?
\end{Ej}

\section{The category over schemes and its fibers}

A \emph{category over $\Sch/\CC$} starts with a category and asks
us to specify a functor from it to $\Sch/\CC$. It need not yet have
pullbacks or satisfy any descent condition. Here define
\[
p:\Mgn\longrightarrow\Sch/\CC,\qquad
p(\xi)=S,\qquad p(u,h)=u.
\]
The identity calculation gives $p(\id_\xi)=\id_S$, and the
composition calculation gives $p((v,k)\circ(u,h))=v\circ u$.
Thus $p$ is a covariant functor. If this is the functor previously
called \texttt{Target}, then $p=\texttt{Target}$: it retains the
target $S$ of the structural morphism $C\to S$. Later, for stable
maps, this word will still mean $S$, not the fixed target variety $X$.
This is Fantechi's Section~3.2 notion of a category over the base
\cite{Fantechi}.

\begin{Def}[The fiber category]\label{def:curve-fiber}
Fix $S$. The category $\Mgn(S)$ has all families over this exact
scheme $S$ as objects. Its arrows are the arrows of the total
category lying over $\id_S$. Thus an arrow is an $S$-isomorphism
$h:C\xrightarrow{\sim}D$ satisfying $h\sigma_a=\tau_a$.
\end{Def}

The word ``fiber'' here means a category, not the geometric fiber
$C_{\bar s}$ of a scheme morphism. Objects in $\Mgn(S)$ can be
entire varying families. For $S=\Spec\CC$, they are individual
stable pointed curves and their marked isomorphisms. The requirement
that arrows lie over the \emph{identity} is important: an automorphism
of $S$ different from the identity would give an arrow in the total
category, but not in this fiber category.

\section{Pullbacks and cartesian arrows}\label{sec:curve-pullback}

\subsection{Construct the changed family}

Suppose we have a family $\xi$ over $S$ and a $\CC$-morphism
$u:T\to S$. We are asked to produce its family over $T$.
Set $C_T=T\times_SC$, let $\pi_T:C_T\to T$ be projection, and
put $\sigma_{a,T}=(\id_T,\sigma_a u)$. Then
\[
\begin{tikzcd}[row sep=large,column sep=large]
C_T=T\times_SC\arrow[r,"q"]\arrow[d,"\pi_T"']&C\arrow[d,"\pi"]\\
T\arrow[r,"u"']&S
\end{tikzcd}
\]
is a cartesian square and $q\sigma_{a,T}=\sigma_a u$.

Properness, flatness, and finite presentation survive base change.
So do the smooth-locus and disjointness conditions on the sections.
The geometric fibers of $\pi_T$ are geometric base changes of those
of $\pi$; connectedness, nodality, genus, and the component stability
test are unchanged by extension of algebraically closed fields.
Consequently $u^*\xi=(C_T\to T,\sigma_{\bullet,T})$ is an allowed
object. The pair $(u,q):u^*\xi\to\xi$ is an allowed arrow.
These base-change statements are standard family-of-curves facts;
see \cite[Section 6.3.4, pp.~231--232]{Alper}.

\begin{Ex}[Pulling back the genus-two family]
For the constant family $C_0\times S\to S$, changing the base to
$T$ gives $C_0\times T\to T$. Pullback to a point of $S$ selects
a fiber, while pullback to an infinitesimal scheme such as
$\Spec\CC[\epsilon]/(\epsilon^2)$ also remembers infinitesimal
parameter information. Pullback is thus more than substituting
complex numbers into an equation.
\end{Ex}

\subsection{The categorical universal property}

The word ``cartesian'' has a second, related usage, for an arrow
in a category over a base. Its definition starts with an arrow
$q:\xi_T\to\xi$ over $u:T\to S$. The question is whether all
arrows whose base maps factor through $u$ factor uniquely through
$q$ in the corresponding way.

\begin{Def}[Cartesian arrow]\label{def:cartesian-arrow}
An arrow $q:\xi_T\to\xi$ over $u:T\to S$ is cartesian if for
every object $\zeta$ over $R$, base map $v:R\to T$, and arrow
$a:\zeta\to\xi$ over $uv$, there is a unique arrow
$b:\zeta\to\xi_T$ over $v$ with $q b=a$.
\[
\begin{tikzcd}[row sep=large,column sep=large]
\zeta\arrow[rr,"a"]\arrow[dr,dashed,"b"']&&\xi\\
&\xi_T\arrow[ur,"q"']&
\end{tikzcd}
\qquad
\begin{tikzcd}[row sep=large,column sep=large]
R\arrow[rr,"uv"]\arrow[dr,"v"']&&S\\
&T\arrow[ur,"u"']&
\end{tikzcd}
\]
The base factorization $v$ is specified; it is not something to be
chosen as part of the asserted uniqueness.
\end{Def}

Let us verify this for the pullback just constructed. Write the
total space of $\zeta$ as $E$, its projection as $\lambda:E\to R$,
and the total-space map of $a$ as $a_C:E\to C$. There is exactly
one scheme map
\[
           b_C=(v\lambda,a_C):E\longrightarrow T\times_SC
\]
with those two projections, by the scheme fiber-product property.
It respects the marked sections because both of its projections
do. It remains to check that it defines a \emph{moduli arrow}:
since $a$ is an arrow, $E\simeq R\times_SC$; on the other hand,
\[
                R\times_T(T\times_SC)\simeq R\times_SC.
\]
Thus $E\to R\times_TC_T$ is an isomorphism, exactly the required
cartesian-square condition for $(v,b_C)$. Uniqueness as a scheme
map implies uniqueness as a moduli arrow. This proves the full
universal property, including the requirement that the candidate
factorization belongs to our category.

\begin{Ej}\label{ex:cartesian-uniqueness}
If two pullback objects $\xi_T,\xi'_T$ with arrows to $\xi$ satisfy
Definition~\ref{def:cartesian-arrow}, show that there is a unique
isomorphism between them commuting with the arrows to $\xi$.
Why is it incorrect to omit ``commuting with the arrows'' when
asserting uniqueness for a genus-two curve with an involution?
\end{Ej}

\section{Why the category is fibered in groupoids}

A category over schemes is \emph{fibered} if an object over $S$
has a cartesian pullback arrow over every map $T\to S$. We have
constructed one and checked the universal property. A
\emph{groupoid} is a category in which every arrow is invertible.
In $\Mgn(S)$, the cartesian-square requirement over $\id_S$ says
precisely that $h:C\to D$ is an isomorphism. Its inverse is over
$S$ and satisfies $h^{-1}\tau_a=\sigma_a$, so it too is a fiber
arrow. Therefore all fibers are groupoids.

We have proved that $p:\Mgn\to\Sch/\CC$ is a \emph{category
fibered in groupoids}. In fact every arrow in this total category
is cartesian, since its definition identifies its source with the
appropriate pullback. Thus it also satisfies Fantechi's
Definition~3.1 in Section~3.5 \cite{Fantechi}: liftings with a
fixed target and base map are unique up to a unique compatible
isomorphism. Chosen iterated pullbacks need not be literally equal:
$(uv)^*\xi$ and $v^*u^*\xi$ have their canonical isomorphism.
All overlap equations below use those identifications, as Fantechi
explains in Sections~4.1 and~6.3.

\MCcheckpoint{Pullback explains how to move a known family to a new
base. A stack needs the reverse operation too: recover global data
from compatible local data. Neither part of that gluing assertion
has been proved by merely checking the category axioms.}

\section{What the two stack conditions ask}\label{sec:fantechi}

Fix an \'etale cover $\{S_i\to S\}$. Write
$S_{ij}=S_i\times_SS_j$ and
$S_{ijk}=S_i\times_SS_j\times_SS_k$. These are the places where
two or three descriptions of the same base point can be compared.
For an open cover they are intersections. For a noninjective
cover they are actual fiber products, with several possible
comparisons even when there is only one index $i$.

\begin{Def}[Fantechi's stack conditions]\label{def:stack}
Start with a category fibered in groupoids $\mathcal F$ over our site.
It is a stack if the following two requirements hold.
\begin{enumerate}
\item \emph{Isomorphisms are a sheaf.} Given two objects $\xi,\eta$
already over $S$, and local isomorphisms
$a_i:\xi|_{S_i}\to\eta|_{S_i}$ agreeing over every $S_{ij}$,
there is exactly one global isomorphism $a:\xi\to\eta$ whose
restrictions are the $a_i$.
\item \emph{Every object descent datum is effective.} Given objects
$\xi_i$ over $S_i$ and isomorphisms
$\alpha_{ij}:\xi_i|_{S_{ij}}\to\xi_j|_{S_{ij}}$ satisfying
\begin{equation}\label{eq:cocycle}
\alpha_{jk}\circ\alpha_{ij}=\alpha_{ik}\quad\hbox{over }S_{ijk},
\end{equation}
there is an object $\xi$ over $S$ and isomorphisms
$b_i:\xi|_{S_i}\to\xi_i$ such that
\begin{equation}\label{eq:effective}
             \alpha_{ij}=b_j\circ b_i^{-1}\quad\hbox{over }S_{ij}.
\end{equation}
\end{enumerate}
\end{Def}

This is Fantechi \cite[Sections 4.2--4.3, Definitions 4.1--4.3]{Fantechi}.
In the requested local \path{StacksForEverybody.tex}, the descent
definition is labelled \texttt{def:4.1} (lines 207--211), the
isomorphism-sheaf condition is \texttt{def:4.2} (lines 219--221),
and the stack definition is \texttt{def:4.3} (lines 225--227).
The source-PDF markers place them on pp.~6, 6, and 7, respectively.

The first condition begins with global objects and asks us to glue
\emph{an arrow}. The second begins with local objects and asks us
to build \emph{a global object}. A list of isomorphism classes of
local objects is insufficient for the second: the actual transition
isomorphisms and their cocycle are part of the input. On a triple
overlap, going from description $i$ to $j$ and then to $k$ must
agree with going directly from $i$ to $k$.

For fixed $\xi,\eta$ over $S$, the sheaf in the first condition is
the assignment
\[
(T\xrightarrow{v}S)\longmapsto
\Iso_T(v^*\xi,v^*\eta),
\]
with restriction maps given by pullback. Each value is a \emph{set
of isomorphisms}, not a set of curves. The sheaf assertion means
existence and uniqueness of gluing for this assignment on covers
of every $T$. It makes no assertion here that this sheaf is a
scheme; representability will be a further geometric property.

\section{First stack condition: isomorphisms form a sheaf}\label{sec:curve-isom}

Let $\xi=(C\xrightarrow{\pi}S,\sigma_\bullet)$ and
$\eta=(D\xrightarrow{\rho}S,\tau_\bullet)$ be stable pointed families. Suppose
$a_i:C_i=C\times_SS_i\xrightarrow{\sim}D_i=D\times_SS_i$
are $S_i$-isomorphisms preserving all labels, with
$a_i|_{S_{ij}}=a_j|_{S_{ij}}$. We must construct exactly one
$S$-isomorphism with these restrictions.

\paragraph{The theorem that glues the underlying maps.}
For schemes $Y,Z$ and an \'etale cover $\{Y_i\to Y\}$, compatible
maps $Y_i\to Z$ glue to a unique map $Y\to Z$. More generally this
holds for fpqc covers: Vistoli \cite[Proposition 4.31, pp.~91--92]{Vistoli}
uses the representable-sheaf theorem, and
\cite[Lemma 35.13.7, \MCtag{023Q}]{Stacks} states that theorem
precisely. On affine pieces it is faithfully flat descent of ring
homomorphisms, including the equalizer
$A\to B\rightrightarrows B\otimes_AB$; the theorem extends the
affine calculation to scheme morphisms. This is our external
morphism-descent input, not a consequence of the stack we are
trying to prove exists.

\paragraph{Apply it to the correct covering.}
The maps $C_i\to C$ are an \'etale cover, by base change from
$S_i\to S$. Compose $a_i$ with $D_i\to D$ to get maps
$C_i\to D$. Their overlaps agree by assumption, so the theorem
gives a unique scheme map $a:C\to D$. The two maps from $C$ to
$S$ given by $\rho a$ and $\pi$ agree after this cover; uniqueness
in the same theorem gives $\rho a=\pi$. Thus $a$ is over $S$,
and its actual pullback to $S_i$ is $a_i$.

\paragraph{Recover invertibility.}
The inverse maps $a_i^{-1}$ also agree on overlaps. Applying the
theorem to the cover $\{D_i\to D\}$ gives $b:D\to C$.
Both $ba$ and $\id_C$ restrict to $\id_{C_i}$, hence are equal;
likewise $ab=\id_D$. Consequently $a$ is an isomorphism. We have
not inferred invertibility merely from a bijection on fibers.

\paragraph{Recover the markings and uniqueness.}
For each marking index $r$, the maps $a\sigma_r,\tau_r:S\to D$
agree on all $S_i$, so they are equal. Hence the global
isomorphism is an arrow in $\Mgn(S)$. Any other such arrow with
the same local restrictions is equal as a scheme map to this
one. This proves existence, the required properties, and
uniqueness. The argument works over any test scheme $T\to S$,
so it proves the sheaf condition in full.

\begin{Ex}[Compatible genus-two isomorphisms]
Let $S=\mathbb A^1$, with $U=S\setminus\{0\}$ and
$V=S\setminus\{1\}$. Take $C=D=C_0\times S$. The local maps
$\iota\times\id_U$ and $\iota\times\id_V$ agree on
$U\cap V$ and glue to $\iota\times\id_S$. If instead we take
the identity on $U$ and $\iota$ on $V$, they disagree on the
overlap and no global map has those restrictions. The sheaf
condition never promises to glue incompatible arrows. This
example illustrates the general proof; it does not prove the
condition for other families or other covers.
\end{Ex}

\begin{Ej}\label{ex:iso-sheaf}
Why is it not enough to glue just the underlying maps and say
``therefore they are isomorphisms''? Explain both a way to finish
using their inverses and the separate check required by markings.
\end{Ej}

% Descent tools used before the effective-descent proof in Chapter 1.
\section{The geometric tool that makes object descent possible}
\label{desc:tools}

We now need a theorem that produces a scheme from local schemes.
Before applying it to our curves, we explain its input: a line bundle
that supplies compatible projective coordinates. Merely knowing that
every local curve admits \emph{some} projective embedding does not supply
compatible embeddings on overlaps. The particular line bundle below is
determined by the pointed curve itself, so every transition isomorphism
automatically acts on it.

\begin{Def}[The line bundle supplied by a stable pointed family]
\label{desc:canonical-bundle}
Start with the family $(p:C\to S,\sigma_1,\ldots,\sigma_n)$ of
Definition~\ref{def:stable-family}. Each marking has an image
$D_a=\sigma_a(S)$ in $C$. We form
\[
 D=D_1+\cdots+D_n,
 \qquad L_C=\omega_{C/S}(D)
       =\omega_{C/S}\otimes\mathcal O_C(D).
\]
Here $\omega_{C/S}$ is the \emph{relative dualizing sheaf}: the line
bundle playing the role of relative differential forms for a family
whose fibers are allowed to have nodes. The notation $(D)$ means
tensoring with the line bundle associated to the indicated divisor.
\end{Def}

This definition uses two results about families of schemes, which we
state as external inputs. A section through the smooth locus of a
separated family of curves is a relative effective Cartier divisor.
Concretely, its image is locally cut out by one equation that is a
non-zero-divisor, and this description remains a Cartier divisor after
base change. One can check this in a relative smooth coordinate, where
the section has equation $z=0$; the general statement follows by
\'etale localization. See Stacks Project, Section 44.3,
Lemma 44.3.1, and Section 31.23, Lemma 31.23.8
\cite{Stacks}. Thus all the $\mathcal O_C(D_a)$ above are defined and
commute with base change.

The second input is the relative duality theorem in its curve form:
if $p:C\to S$ is proper, flat, and finitely presented, with
Gorenstein geometric fibers of pure dimension one, then
$\omega_{C/S}$ is invertible and, for every $b:T\to S$, there is a
canonical isomorphism
\begin{equation}\label{desc:dualizing-basechange}
 \operatorname{pr}_C^*\omega_{C/S}
       \simeq \omega_{C\times_S T/T}.
\end{equation}
These isomorphisms respect further base change and isomorphisms of
families. This is a theorem about proper flat families; it does not
require us to know that our moduli category is algebraic. A precise
curve formulation is Stacks Project, \MCtag{0E6R},
Lemma 109.19.3, with its duality references in Section 109.19
\cite{Stacks}. We use this duality statement as a black box.

Why does it apply? A smooth curve has regular local rings; a node is
locally a hypersurface singularity, with completed local ring
$\overline{k}[[x,y]]/(xy)$. These rings are Gorenstein: their dualizing
modules have rank one and are free. Consequently the geometric fibers
in our definition satisfy the theorem. On the smooth locus,
$\omega_{C/S}$ agrees with relative differentials. On a single nodal
fiber its sections can be described on the normalization as
differentials with at most simple poles over each node, with the two
residues summing to zero. This explains why ordinary regular
differentials on a smooth curve need a replacement at a node.

\begin{Lem}[Functorial relative ampleness]
\label{desc:ample-canonical}
For a stable pointed family, $L_C=\omega_{C/S}(D)$ is relatively ample
over $S$. It commutes canonically with arbitrary base change. A pointed
isomorphism $a:C\xrightarrow{\sim}C'$ over $S$ therefore gives a
canonical isomorphism $a^*L_{C'}\simeq L_C$, compatible with
composition of pointed isomorphisms.
\end{Lem}

\noindent\emph{Verification.}
The base-change assertion follows from
\eqref{desc:dualizing-basechange} and from the relative Cartier divisor
property of the sections. For ampleness, restrict to a geometric fiber
and normalize any one of its irreducible components. If that
normalization has genus $h$ and $r$ distinguished points, counting the
branches of nodes and the markings, then
\begin{equation}\label{eq:component-degree}
 \deg L_C|_{\widetilde E}=2h-2+r>0
\end{equation}
by stability. Positivity on every component implies
ampleness on the fiber. The precise curve criterion is Stacks Project,
\MCtag{0B5Y}, Lemma 33.44.15; degrees are unchanged on normalization
by Lemma 33.44.4. For a proper family, ampleness on a fiber implies
relative ampleness on a neighborhood of its base point
(\MCtag{0D2S}, Lemma 37.50.3). These neighborhoods cover $S$, which
proves relative ampleness. The final assertion follows because an
isomorphism carries each labelled section to the corresponding labelled
section, and the duality identifications are functorial.\hfill$\square$

In ordinary language, every stable pointed family comes with the
\emph{same rule} for choosing an ample line bundle. We have not added a
line bundle to the objects of $\Mgn$: we have constructed one from data
already present. For a smooth unpointed genus-two curve $C$, this rule
is simply $L_C=\omega_C$, of degree $2$. For a genus-two curve made from
two smooth elliptic curves joined at one node, its degree on each
component is $1$. The positivity check works on both components even
though the curve is singular.

\begin{Th}[Effective descent with a compatible ample line bundle]
\label{desc:polarized-theorem}
Let $\{S_i\to S\}$ be an \textup{\'etale} cover of a scheme.
Suppose that for every $i$ we are given a proper, flat, finitely
presented morphism $p_i:Y_i\to S_i$ and a relatively ample invertible
sheaf $L_i$ on $Y_i$. Suppose also that there are isomorphisms
\[
 a_{ij}:Y_i|_{S_{ij}}\xrightarrow{\sim}Y_j|_{S_{ij}},
 \qquad
 \lambda_{ij}:a_{ij}^*(L_j|_{S_{ij}})
                   \xrightarrow{\sim}L_i|_{S_{ij}},
\]
where $S_{ij}=S_i\times_S S_j$. Assume the scheme isomorphisms satisfy
$a_{ik}=a_{jk}\circ a_{ij}$ on $S_{ijk}$, and the line-bundle
isomorphisms satisfy
\begin{equation}\label{desc:bundle-cocycle}
 \lambda_{ik}=\lambda_{ij}\circ a_{ij}^*\lambda_{jk}
 \quad\text{on }S_{ijk},
\end{equation}
using the canonical identifications of iterated pullbacks.
Then there are a proper, flat, finitely presented scheme $Y\to S$, a
relatively ample invertible sheaf $L$ on $Y$, and isomorphisms
$(Y,L)|_{S_i}\simeq(Y_i,L_i)$ recovering the prescribed descent data.
The same statement holds for fpqc covers.
\end{Th}

\noindent\textbf{Source and proof strategy.}
The scheme-and-line-bundle effectivity theorem is SGA~1,
Expos\'e VIII, Proposition 7.8, p.~174 of the recomposed edition
\cite{SGA1}. Its argument descends a graded algebra and uses
relative $\Proj$. Vistoli's later exposition \cite{Vistoli},
Section 4.3.3, Theorem 4.38 and proof, pp.~96--103, gives the detailed
version for a functorial ample bundle on a class of proper flat
finitely presented families. His theorem applies this method to a
whole moduli problem; the statement above isolates the descent datum
that we need. Neither source says that arbitrary projective schemes
descend effectively without compatible polarization data.

\noindent\textbf{Our guided explanation of the construction.}
The input is local spaces together with instructions for identifying
them and their line bundles on overlaps. The output must be one
scheme over $S$, with identifications that reproduce precisely those
instructions. We describe the \emph{\'etale} case used in these notes.

First work over an affine open of $S$. Refine the cover by affine open
subschemes of its members. Because \'etale maps are open and $S$ is
quasi-compact, finitely many of their images cover $S$. Their disjoint
union $U$ is affine, and $U\to S$ is faithfully flat and
quasi-compact. The map need not be finite: it is the \emph{list of
affine covering schemes} that is finite. All the original data restrict
to $U$ and $U\times_S U$. After carrying out descent there, the
uniqueness of morphism descent identifies the answer with the data on
every original $S_i$. Answers over different affine opens of $S$ glue
by ordinary Zariski gluing. This reduction justifies using one uniform
power of a line bundle in the following steps.

\begin{enumerate}
\item \emph{Turn the local families into closed subschemes of projective
bundles.}
Write $p_U:Y_U\to U$ and $L_U$ for the disjoint union of the refined
data. Relative Serre vanishing and the relative very-ampleness theorem
allow us to choose $m>0$ such that $L_U^{\otimes m}$ is relatively very
ample and $H^j(Y_u,L_u^{\otimes m})=0$ for every $j>0$ on every fiber. Cohomology and
base change then makes
\[
 E_U=(p_U)_*L_U^{\otimes m}
\]
a finite locally free sheaf, with formation commuting with base change.
These are substantial scheme-theoretic inputs, not consequences of the
stack axioms. The cohomology statement, including arbitrary bases, is
Vistoli, Proposition 4.37, pp.~95--96; the passage to a sufficiently
large power occurs in the proof of Theorem 4.38. The evaluation map
$p_U^*E_U\to L_U^{\otimes m}$ gives a closed immersion
\[
 Y_U\hookrightarrow\PP_U(E_U)
           :=\Proj_U\bigl(\operatorname{Sym}^{\bullet}E_U\bigr).
\]
Here relative $\Proj$ is the construction that turns a graded algebra
over a base into a projective scheme over that base.

\item \emph{Descend the projective coordinates.}
The isomorphisms $a_{ij}$ and $\lambda_{ij}$ induce compatible
isomorphisms between the pullbacks of $E_U$ on $U\times_S U$.
Equation~\eqref{desc:bundle-cocycle} gives their cocycle condition.
Faithfully flat descent for quasi-coherent modules now constructs a
sheaf $E$ on $S$ whose pullback is $E_U$. Local freeness and finite rank
can be checked faithfully flat locally, so $E$ is a vector bundle.
We use Vistoli, Theorem 4.23, pp.~85--90, for this descent theorem.
It is the algebraic input that replaces elementary gluing of vector
bundles on open subsets.

Thus the projective bundles are all pullbacks of the single scheme
$\PP_S(E)$. Functoriality of evaluation ensures that the given
isomorphisms identify the embedded copies of $Y_U$ in this common
projective bundle. This is the step for which compatible line bundles
were essential.

\item \emph{Descend the equations.}
Let $\mathcal I_U\subset\mathcal O_{\PP_U(E_U)}$ be the ideal sheaf
of the embedded $Y_U$. The compatibility just checked gives this ideal
sheaf a descent datum over the fpqc cover
$\PP_U(E_U)\to\PP_S(E)$. Quasi-coherent descent constructs an ideal
sheaf $\mathcal I\subset\mathcal O_{\PP_S(E)}$; both its inclusion
and its closure under multiplication can be checked after the cover.
Let $Y\subset\PP_S(E)$ be the closed subscheme cut out by
$\mathcal I$. Formation of this closed subscheme commutes with flat
base change, so $Y\times_S U\simeq Y_U$ with the \emph{specified}
transition identifications. This is the projective-bundle form of the
graded-algebra/$\Proj$ construction in SGA~1.

\item \emph{Recover the line bundle and check the family.}
Now that $Y$ exists, the $L_i$ are line bundles on a cover of $Y$ with
compatible transition maps. Quasi-coherent descent constructs $L$;
invertibility is checked on that cover. Properness, flatness, and finite
presentation of $Y\to S$ follow from their fpqc locality on the base
(Stacks Project, \MCtag{02L1}, \MCtag{02L2}, \MCtag{02L0},
Lemmas 35.23.16, 35.23.17, 35.23.15). Relative ampleness of $L$
can be checked on geometric fibers and then near each base point
using the proper-family ampleness theorem above.
\end{enumerate}

This explains why the theorem returns a \emph{scheme}. It does not
begin by assuming that the sheaf quotient of the $Y_i$ is a scheme.
The vector bundle and the descended equations construct that scheme.
The theorem has not yet supplied markings or a map to a target; we
will descend those additional morphisms after applying it.

\begin{Ej}\label{desc:exercise-polarization}
Suppose the local pointed families have transition isomorphisms
$a_{ij}$ satisfying the cocycle condition. Explain why choosing
$L_i=\omega_{C_i/S_i}(\sum_a\sigma_{a,i})$ supplies the
$\lambda_{ij}$ and equation~\eqref{desc:bundle-cocycle} without an
additional choice. Why would choosing an unrelated ample line bundle
on each $C_i$ not by itself do this?
\end{Ej}


\section{Second stack condition: descent data are effective}\label{sec:curve-descent}

The input now differs from that of the preceding sheaf proof.
There need not be a known global curve. We are given an \'etale
cover $\{S_i\to S\}$, families
\[
 (C_i\xrightarrow{\pi_i}S_i,\sigma_{1,i},\ldots,\sigma_{n,i}),
\]
and isomorphisms
\[
\alpha_{ij}:C_i\times_{S_i}S_{ij}
   \xrightarrow{\sim}C_j\times_{S_j}S_{ij}
\]
over $S_{ij}$ taking every $\sigma_{a,i}$ to $\sigma_{a,j}$.
On $S_{ijk}$, we are given the equality
$\alpha_{jk}\alpha_{ij}=\alpha_{ik}$ of scheme maps.
That last requirement means the following triangle commutes:
\[
\begin{tikzcd}[row sep=large,column sep=large]
C_i|_{S_{ijk}}\arrow[r,"\alpha_{ij}"]\arrow[dr,"\alpha_{ik}"']
 &C_j|_{S_{ijk}}\arrow[d,"\alpha_{jk}"]\\
 &C_k|_{S_{ijk}}.
\end{tikzcd}
\]
All three vertices are curves over the \emph{same} triple overlap.
For a nontrivial \'etale cover, one cannot replace these fiber
products by informal intersections.

\paragraph{Step 1: provide compatible ample bundles.}
On $C_i$ form
$L_i=\omega_{C_i/S_i}(\sum_{a=1}^n\sigma_{a,i})$.
The preceding descent discussion establishes that it is invertible
and relatively ample, commutes with base change, and is carried to
the corresponding bundle by an isomorphism of pointed families.
In particular $\alpha_{ij}^*L_j\simeq L_i$ on the overlap.
These are the canonical pullback identifications, so their
composites on triple overlaps agree. There is no arbitrary
scalar choice of transition maps on the line bundles.

\paragraph{Step 2: construct the global scheme.}
Apply the polarized scheme-descent theorem stated above to
$(C_i,L_i,\alpha_{ij})$. All its hypotheses have just been
checked: the cover is faithfully flat locally of finite
presentation after the stated refinements, each $C_i\to S_i$
is proper, flat, and finitely presented, the bundles are relatively
ample, and both the scheme and bundle transition maps satisfy
their cocycles. We obtain a scheme $C\to S$, a relatively ample
bundle $L$, and isomorphisms
\[
                       b_i:C\times_SS_i\xrightarrow{\sim}C_i
\]
whose transition maps are the given $\alpha_{ij}$. These
identifications are part of the output: effectiveness asks for
them, not just for some global curve with isomorphic fibers.

Properness, finite presentation, and flatness of $C\to S$ can
be checked after the faithfully flat cover, so they hold because
they hold for every $\pi_i$. This uses the descent-of-properties
theorem \cite[Proposition 1.15, pp.~11--12]{Vistoli}, or precisely
\MCtag{02L0}, \MCtag{02L1}, and \MCtag{02L2} in \cite{Stacks}.

\paragraph{Step 3: descend the sections.}
The local sections, transported by $b_i^{-1}$ and then projected
to $C$, give maps $s_{a,i}:S_i\to C$. They agree on $S_{ij}$
because $\alpha_{ij}$ preserves marking $a$. The morphism-sheaf
theorem gives a unique $\sigma_a:S\to C$ with these restrictions.
The equation $\pi\sigma_a=\id_S$ holds after the cover, hence
globally. Thus we have recovered genuine sections, including
over nonreduced parameter schemes.

Their images lie in the smooth locus because this assertion is
local after \'etale base change. They are disjoint: a point in
an intersection would lift after the surjective cover to an
intersection of the corresponding local sections, which does
not exist. Equivalently their fiber product is empty after a
faithfully flat cover, so it was empty to start with.

\paragraph{Step 4: check the geometric fibers and stability.}
Take any geometric point $\bar s:\Spec\Omega\to S$. After
passing to an algebraically closed extension if necessary, it
lifts to some $S_i$. The geometric fiber of $C$ at that lifted
point is, by $b_i$, the corresponding fiber of $C_i$. It is
connected, reduced, nodal, has genus $g$, has distinct smooth
markings, and meets the component stability inequalities.
These geometric assertions are unchanged by algebraically
closed field extension; hence they hold for $C_{\bar s}$.
We have verified every part of Definition~\ref{def:stable-family}.
The recovered $L$ is also canonically the bundle
$\omega_{C/S}(\sum\sigma_a)$, since that identification holds
locally and the compatible bundle maps descend.

\paragraph{Step 5: check the prescribed descent datum.}
By construction $\alpha_{ij}=b_jb_i^{-1}$. The $b_i$ preserve
all recovered sections, so they are isomorphisms in the fiber
categories, not just isomorphisms of unmarked schemes. We have
therefore produced exactly the object and identifications
required by Fantechi's Definition~4.1.

If another family $C'\to S$ with identifications $b_i'$ has
the same descent datum, the local maps $b_i^{-1}b_i'$ agree
on overlaps. The first stack condition glues them to a unique
global isomorphism compatible with all identifications. This
is the appropriate uniqueness statement; the object need not
have a trivial automorphism group.

\subsection{A nontrivial genus-two descent datum}\label{sec:genus-two-twist}

Let $S=\Spec\CC[t,t^{-1}]$ and let
$U=\Spec\CC[u,u^{-1}]\to S$ be $t=u^2$. Over $U$ take
the constant family $C_0\times U$. The overlap $U\times_SU$
has two components, specified by $u_2=u_1$ and $u_2=-u_1$.
Define the transition map to be the identity on the first
component and $\iota$ on the second. In affine coordinates it is
\[
                  (x,y_1)\longmapsto(x,y_2),
                   \qquad y_2=(u_1/u_2)y_1.
\]
On a triple overlap the factors multiply as
$(u_2/u_3)(u_1/u_2)=u_1/u_3$, so the cocycle holds.
The data are a family and a rule comparing its two possible
square-root descriptions, even though the cover has only one
member.

Here the descended family can be written explicitly as the
smooth projective double cover with affine equation
\[
                            z^2=t(x^6-1).
\]
To specify the completion over the whole base, use the section
$t(X^6-Z^6)$ of $\mathcal O_{\PP^1_S}(6)$ and the usual
double-cover algebra
$\mathcal O_{\PP^1_S}\oplus\mathcal O_{\PP^1_S}(-3)$,
whose multiplication is determined by that section. Its relative
spectrum is finite over $\PP^1_S$, hence proper over $S$.
After the \'etale cover it is $C_0\times U$, so it is flat,
smooth of relative dimension one, and genus two. The local
identification is $y=z/u$; these identifications have exactly
the stated transition maps. The section $(x,z)=(1,0)$ can be
included if we want the one-pointed version.

Pullback along $T\to S$ replaces $t$ by its invertible pullback
function on $T$. In particular the original square-root cover
makes the family constant. The example shows why transition
maps are real information: locally constant descriptions can
carry a nontrivial gluing rule. Our general descent theorem,
not this equation, was what proved effectiveness for every
stable-curve datum.

\begin{Ej}\label{ex:twisted-datum}
For this one-map cover, explain why the overlap still contains
more than the diagonal. Check the cocycle when the first
comparison and the second comparison are both sign changes.
Compute the transition on the local expression $y=z/u$.
\end{Ej}

\begin{Th}[The stable pointed-curve stack]\label{th:curve-stack}
For $g,n\geq0$ with $2g-2+n>0$, the category of stable
$n$-pointed genus-$g$ curves defined above, with its projection
to $\Sch/\CC$, is a stack for the \'etale topology.
\end{Th}
\begin{proof}
Identities and composition make it a category. The projection
is a functor. Base change gives cartesian liftings and all
fibers are groupoids. Compatible isomorphisms glue uniquely
by Section~\ref{sec:curve-isom}. Every object descent datum
is effective by Section~\ref{sec:curve-descent}. These are
exactly the hypotheses and the two conclusions required by
Definition~\ref{def:stack}.
\end{proof}

\section{Properties}\label{sec:curve-properties}

The stack proof gives an answer to a gluing question: compatible local
families, and compatible local isomorphisms, come from global ones.
We now ask geometric questions about that stack. These require additional
theorems. Throughout this section $g,n\geq 0$, $2g-2+n>0$, and
$\Mgn$ is the stack over $\CC$ just constructed.

\subsection{What the geometric words ask us to construct}

\begin{Def}[Representability and an atlas]
Start with a morphism $F:\mathcal Y\to\mathcal Z$ of stacks. To test
whether it is \emph{representable by algebraic spaces}, choose any scheme
$S$ and any object $z\in\mathcal Z(S)$, equivalently a map
$S\to\mathcal Z$, and form
\[
\begin{tikzcd}
S\times_{\mathcal Z}\mathcal Y \arrow[r]\arrow[d]
  & \mathcal Y\arrow[d,"F"]\\
S\arrow[r,"z"] & \mathcal Z.
\end{tikzcd}
\]
The requirement is that the fiber product be an algebraic space.
An algebraic space is an \'etale sheaf which has a representable,
surjective \'etale presentation
by a scheme; it allows \'etale coordinate charts even when those charts
do not glue to a scheme by open immersions. In the applications below
these particular fiber products will actually be schemes.

If $F$ is representable, saying that it is smooth, \'etale, or
surjective means that the resulting maps of algebraic spaces to $S$
have that property for every such test. An \emph{atlas} is a
representable surjective map $U\to\mathcal Z$ from a scheme, with the
specified smoothness or \'etaleness condition.
\end{Def}

For example, a map $U\to\Mgn$ is a family of stable pointed curves
over $U$. Given another family $C/S$, a point of
$S\times_{\Mgn}U$ consists of a parameter $u$ in $U$ and an
\emph{isomorphism} from the curve selected by $u$ to the corresponding
fiber of $C/S$. The isomorphism is part of the data. The fiber product
does not merely record that two isomorphism classes agree.

\begin{Def}[Algebraic, Artin, and Deligne--Mumford]
Start with a stack $\mathcal Z$ on the \'etale site of schemes over
$\CC$. In these notes, it is an \emph{algebraic stack}, also called
an \emph{Artin stack}, if its diagonal is representable by algebraic
spaces and it has a smooth atlas $U\to\mathcal Z$ from a scheme.
It is a \emph{Deligne--Mumford stack} if it has an \'etale atlas.
Thus every Deligne--Mumford stack is an algebraic stack in our
convention, since \'etale morphisms are smooth.
\end{Def}

The diagonal has a concrete interpretation here. Pulling
$\Delta:\Mgn\to\Mgn\times\Mgn$ back by two families $C/S,D/S$
gives the sheaf $\Iso_S(C,D)$, with the markings understood.
The earlier sheaf proof showed that this functor is a sheaf.
Representability of the diagonal asks for more: this sheaf must
itself be an algebraic space. An atlas asks for still more: enough
families must be parametrized by schemes in a smooth or \'etale way.
For these conventions and their relation, see Alper
\cite[Definitions 4.1.3--4.1.6 and Theorem 4.2.1, pp.~105--106,
118--119]{Alper}. Some older sources, including Deligne--Mumford,
use ``algebraic stack'' for what we now call a Deligne--Mumford stack.

\subsection{Algebraicity and finite type}

\begin{Prop}\label{prop:curves-algebraic}
The stack $\Mgn$ is algebraic and of finite type over $\CC$.
Its diagonal is affine and of finite type. It has a smooth
surjective atlas $H\to\Mgn$, where $H$ is a scheme of finite
type parametrizing suitable embedded stable pointed curves.
\end{Prop}

\noindent\textbf{Original source and proof strategy.}
Deligne--Mumford \cite[Section 1, pp.~77--78, and Proposition (5.1),
p.~104]{DM} construct the unpointed stack using tricanonical Hilbert
schemes. Knudsen \cite[Theorem 2.7, p.~179]{Knudsen} proves the
pointed theorem over $\operatorname{Spec}\mathbb Z$, using the
unpointed result and induction in the number of markings.
We give a direct Hilbert-atlas argument for pointed stable curves.
Its embedding bound is Alper
\cite[Proposition 6.3.23, p.~231]{Alper}; its smooth-atlas proof is
the embedding-lifting argument of
\cite[Theorem 6.4.6, pp.~238--240]{Alper}, applied here directly
 to our stable category. Although that theorem is written for a
larger category of curves, we do not construct that larger stack.
Our parameter scheme retains some extra choices of embedding;
these choices are allowed in an atlas.

\medskip
\noindent\textbf{Our guided proof.}
The data already available are arbitrary stable pointed families
$(\pi:C\to S,\sigma_1,\ldots,\sigma_n)$. We must construct one
scheme containing enough such families and show that its map to
our stack is representable, smooth, and surjective.

\smallskip
\noindent\emph{Step 1: use a canonical line bundle to bound embeddings.}
Put
\[
 L_C=\omega_{C/S}\Bigl(\sum_{a=1}^n\sigma_a\Bigr).
\]
We use the following \emph{external embedding theorem}: for a stable
pointed family, $L_C^{\otimes3}$ is relatively very ample,
$E_C=\pi_*L_C^{\otimes3}$ is locally free of rank $r=5g-5+3n$,
and formation of this vector bundle commutes with arbitrary base
change. These conclusions follow from the fiberwise vanishing and
very ampleness theorem together with cohomology and base change;
see Alper \cite[Exercise 6.3.18 and Proposition 6.3.23,
pp.~230--231]{Alper}. The fiberwise unpointed theorem is proved in
Deligne--Mumford \cite[Theorem (1.2) and its corollary,
pp.~77--78]{DM}. We take the pointed very ampleness result as an
input rather than redo its separation-of-points-and-tangents proof.

The rank and Hilbert polynomial have a useful numerical explanation.
On any geometric fiber, let $\ell=2g-2+n>0$. Then
$\deg L_C=\ell$ and $\deg L_C^{\otimes3}=3\ell$. Vanishing of
$H^1(C,L_C^{\otimes3})$ and Riemann--Roch give
\[
 r=3\ell+1-g=5g-5+3n,
 \qquad P(m)=\chi(C,L_C^{\otimes3m})=3\ell m+1-g.
\]
These numbers depend only on $g,n$, not on the particular curve.
The evaluation of sections gives a closed immersion
$C\hookrightarrow\mathbb P_S(E_C)$, where our convention for
$\mathbb P(E_C)$ parametrizes one-dimensional quotients of $E_C$.
Locally on $S$, choosing a basis of $E_C$ identifies the ambient
bundle with $\mathbb P^{r-1}_S$.

\smallskip
\noindent\emph{Step 2: represent the diagonal.}
Take two stable pointed families $C/S,D/S$. After trivializing
$E_C,E_D$ locally over $S$, their canonical embeddings place them
in the same projective space. An isomorphism $C\to D$ induces an
isomorphism of their canonical line bundles, hence a unique
projective linear transformation of these complete embeddings.
Conversely, a projective transformation carrying the embedded
curve and every marking to the corresponding data restricts to
an isomorphism of pointed families.

Thus $\Iso_S(C,D)$ is locally represented by a closed subscheme
of $\PGL_{r,S}$: the transformation must send one Hilbert point
 to the other, and must send each marked point to its partner.
Equality of Hilbert points is closed because the Hilbert scheme
is separated; the point conditions are closed because projective
space is separated. These are finitely presented conditions, since
the Hilbert scheme and projective space are of finite presentation
over $\CC$. The local
representing schemes glue by descent for affine schemes, so
$\Iso_S(C,D)$ is an affine scheme of finite presentation over $S$.
This proves the claimed representability and finiteness of the
diagonal. It also implies that a map from any scheme to our
stack is representable: its test fiber products are obtained by
base changing that diagonal. Notice how much stronger this is
than merely saying that $\Iso_S(C,D)$ is a sheaf.

\smallskip
\noindent\emph{Step 3: choose a finite type parameter scheme.}
The \emph{Hilbert scheme theorem} is an external input:
closed subschemes of $\mathbb P^{r-1}_S$ flat and finitely presented
over $S$, with fiber Hilbert polynomial $P$, are represented by a
projective scheme $\operatorname{Hilb}^{P}(\mathbb P^{r-1}_{\CC})$.
See Alper \cite[Theorem 2.1.2, p.~48]{Alper}.
Adding $n$ points and requiring each to lie on the universal curve
gives a closed incidence scheme in
\[
 \operatorname{Hilb}^{P}(\mathbb P^{r-1}_{\CC})
                      \times(\mathbb P^{r-1}_{\CC})^n.
\]
Inside it, take the open locus $H$ where the pointed fibers are
stable and $H^1(C_s,\mathcal O_{C_s}(1))=0$.
Openness of stability is Alper
\cite[Proposition 6.3.24, p.~232]{Alper}; the vanishing locus is
open by upper semicontinuity of cohomology. Both are substantial
standard inputs about projective flat families. The Hilbert
polynomial already fixes the arithmetic genus to be $g$.
Since the incidence scheme is of finite type over the noetherian
field $\CC$, its open subscheme $H$ is of finite type.

The universal family defines $a:H\to\Mgn$ by forgetting the
embedding. It is representable by Step 2. It is surjective because
any stable family is, locally on its base, embedded by $L_C^3$
as in Step 1, and these embeddings satisfy the defining conditions
of $H$. We deliberately allow other embeddings in $H$ as well;
being an atlas does not require a unique embedding for each curve.

\smallskip
\noindent\emph{Step 4: why forgetting the embedding is smooth.}
Smoothness says that an embedding can be chosen compatibly when a
family undergoes a small infinitesimal variation. We apply the
\emph{infinitesimal criterion for smoothness} to the representable
map $a$, following Alper
\cite[Theorem 4.7.1, pp.~144--145; proof of Theorem 6.4.6,
pp.~238--240]{Alper}. The finite-presentation hypothesis of the
criterion also holds. For a test family $C/S$, the scheme
$S\times_{\Mgn}H$ is the isomorphism scheme between the two
families over $S\times H$: the pullback of $C/S$ and the pullback
of the universal family over $H$. Step 2 makes this scheme
finitely presented over $S\times H$, and $S\times H$ is finitely
presented over $S$. Thus $a$ is representable and locally of finite
presentation before we use the lifting calculation.

Let $A'\twoheadrightarrow A$ be a small extension of local Artinian
$\CC$-algebras, with residue field $K$ and kernel $J$ killed by the
maximal ideal of $A'$. We are given a stable pointed family $C'/A'$
and an embedding of its restriction $C=C'\times_{A'}A$ in
$\mathbb P^{r-1}_A$, representing a point of $H(A)$. We must lift
that embedding to $C'$:
\[
\begin{tikzcd}
 C\arrow[r,hook]\arrow[d,hook] & \mathbb P^{r-1}_A\arrow[d,hook]\\
 C'\arrow[r,dashed,hook] & \mathbb P^{r-1}_{A'}.
\end{tikzcd}
\]
The embedding on $C$ is given by $M=\mathcal O_C(1)$ and its
$r$ coordinate sections.

A line bundle on $C$ lifts to $C'$ because its lifting obstruction
lies in $H^2(C_K,\mathcal O_{C_K}\otimes_KJ)=0$.
This elementary line-bundle lifting theorem follows from
$1+\mathcal I\to\mathcal O_{C'}^*\to\mathcal O_C^*$ for the
square-zero ideal $\mathcal I$; see Alper
\cite[Proposition C.2.11 and Remark C.2.12, p.~602]{Alper}.
Choose a lift $M'$. The restriction sequence is
\[
 0\longrightarrow M_K\otimes_KJ
   \longrightarrow M'\longrightarrow M\longrightarrow0.
\]
The defining vanishing $H^1(C_K,M_K)=0$ implies that every
coordinate section of $M$ lifts to $M'$. The lifted sections
still generate $M'$ by Nakayama's lemma. They define a morphism
$C'\to\mathbb P^{r-1}_{A'}$ lifting the given embedding.

It is a closed immersion. Indeed it is proper, its only geometric
fiber is a closed immersion and hence it is quasi-finite, so it
is finite. The cokernel of the map from the ambient structure
sheaf to its finite pushforward restricts to zero over $A$ and
therefore vanishes by Nakayama's lemma. The resulting embedded
family is flat over $A'$ because it is the given family $C'$.
Its marked sections were already given and acquire their images
under this embedding. Its special fiber lies in $H$, so the
lift also lies in the open locus $H$. This completes the required
lifting. The infinitesimal criterion now proves that $a$ is smooth.

We have constructed a representable smooth surjective map from
the scheme $H$ to $\Mgn$, and have represented the diagonal.
These are exactly the requirements for algebraicity. Finally,
finite type of $H$ proves that $\Mgn$ is quasi-compact and
locally of finite type over $\CC$, which is the asserted finite
type property.

\begin{Ex}[Genus two and extra embedding choices]
For $(g,n)=(2,0)$ we have $\ell=2$, $r=5$, and $P(m)=6m-1$.
Every stable genus-two curve has a tricanonical embedding in
$\mathbb P^4$. Thus the Hilbert scheme for that one polynomial
already contains embeddings of every stable genus-two curve.
The atlas retains choices of an embedding and coordinates;
$\overline{\mathcal M}_2$ retains the abstract curve and its
isomorphisms. Extra parameters in an atlas are harmless: a
smooth atlas need not have relative dimension zero.
\end{Ex}

\subsection{Why the diagonal is unramified}

\begin{Prop}\label{prop:curves-dm}
Under the same hypotheses, $\Mgn$ is Deligne--Mumford.
Every geometric stable pointed curve has a finite reduced
automorphism group scheme.
\end{Prop}

\noindent\textbf{Original source and proof strategy.}
Deligne--Mumford \cite[Lemma (1.4), pp.~80--81, Theorem (1.11),
p.~84, and Proposition (5.1), p.~104]{DM} eliminate infinitesimal
automorphisms and construct the unpointed Deligne--Mumford stack.
Knudsen's pointed theorem is \cite[Theorem 2.7, p.~179]{Knudsen}.
The pointed calculation is explained by Alper
\cite[Proposition 6.3.26, pp.~232--235, and Theorem 6.4.16,
p.~241]{Alper}. The key geometric fact is that every component of
the normalization has enough special points to eliminate regular
vector fields fixing those points.

\medskip
\noindent\textbf{Our guided proof.}
We already have algebraicity. The additional construction we need
is an \'etale atlas. We use the following \emph{Deligne--Mumford
criterion} as an external theorem: an algebraic stack has an
\'etale atlas if and only if its diagonal is unramified, equivalently
if every geometric stabilizer is discrete and reduced. In our
finite type situation these stabilizers will be finite.
See Alper \cite[Theorem 4.6.4, p.~143]{Alper}.

Fix an algebraically closed extension $K/\CC$ and a stable pointed
curve $(C,p_1,\ldots,p_n)$ over $K$. An \emph{infinitesimal
automorphism} starts with the constant family over
$K[\epsilon]/(\epsilon^2)$ and is an automorphism that restricts
to the identity when $\epsilon=0$. Locally on functions it has
the form $a\mapsto a+\epsilon D(a)$, where the multiplication
rule forces $D$ to be a derivation. Fixing a marked point forces
the vector field $D$ to vanish there. Thus the tangent space of
the automorphism scheme at its identity is
\[
 H^0\bigl(C,\mathcal Hom(\Omega_C,\mathcal O_C(-\sum p_a))\bigr).
\]

Let $\nu:\widetilde C\to C$ be the normalization. On a nodal
curve a regular vector field lifts to the normalization and
vanishes at each branch over every node. This can be checked
locally in $K[[x,y]]/(xy)$ by applying a derivation to $xy=0$;
see Alper \cite[Exercise 6.2.22, p.~221]{Alper}.
Write $C_v$ for the smooth connected components of $\widetilde C$,
$h_v$ for their genera, and $B_v$ for the marked points and
branches of nodes on $C_v$. If $m_v=\deg B_v$, stability says
\[
                  2h_v-2+m_v>0.
\]
Consequently
\[
 \deg T_{C_v}(-B_v)=2-2h_v-m_v<0.
\]
A line bundle of negative degree on a smooth projective connected
curve has no nonzero section. Every component vector field
therefore vanishes, so the infinitesimal automorphism vanishes.

The automorphism scheme is an affine group scheme of finite type
by the preceding diagonal calculation. Its zero tangent space at
the identity implies that the identity is an isolated reduced
point: in its noetherian local ring the equality
$\mathfrak m/\mathfrak m^2=0$ implies $\mathfrak m=0$ by
Nakayama's lemma. Translation gives the same conclusion at every
geometric point. A zero-dimensional reduced scheme of finite
type over $K$ has finitely many points. Thus the stabilizer is
finite and reduced, and the Deligne--Mumford criterion applies.

In ordinary language, there can be discrete symmetries of the
curve, but there is no continuously varying infinitesimal symmetry.
For a genus-two hyperelliptic curve, the involution remains an
automorphism; it does not prevent the stack from being
Deligne--Mumford. It does prevent us from identifying that moduli
problem with a scheme which has the same fiber groupoids.

\begin{Ej}\label{ex:curve-dm}
Check the three cases $h_v=0$, $h_v=1$, and $h_v\geq2$ in the
inequality above. Why would a smooth unmarked elliptic curve fail
this argument? Explain why ``finite automorphism groups'' and
``trivial automorphism groups'' are different requirements.
\end{Ej}

\subsection{Properness and the meaning of a limit}

\begin{Prop}\label{prop:curves-proper}
The Deligne--Mumford stack $\Mgn$ is separated and proper over
$\CC$.
\end{Prop}

\noindent\textbf{Original source and proof strategy.}
Deligne--Mumford \cite[Section 2 and Theorem (5.2),
pp.~104--105]{DM} prove the unpointed properness theorem using
stable reduction and uniqueness of stable limits.
Knudsen \cite[Theorem 2.7, p.~179]{Knudsen} proves the pointed
version. A later detailed explanation in characteristic zero is
Alper \cite[Theorem 6.5.1 and Section 6.5.1, pp.~242--247;
Proposition 6.5.20 and Theorem 6.5.23, pp.~251--252]{Alper}.
Alper's written uniqueness proof is unpointed; the marked
extension is assigned in Exercise 6.5.22. We take the full pointed
existence and uniqueness statements from the established pointed
theorem as external inputs. We do not claim to have supplied a
new proof of stable reduction here.

\medskip
\noindent\textbf{Our guided proof.}
For an algebraic stack, \emph{separated} means that its diagonal
is proper. In terms of curves this is a statement about extending
\emph{isomorphisms} between families. \emph{Proper over $\CC$}
means separated, of finite type, and universally closed over
$\CC$. The last condition is a completeness condition: limits
cannot escape after any change of base. To test it effectively we
use the stack valuative criterion
\cite[Theorem 4.8.7, pp.~149--150]{Alper}.

Here is the data in that criterion. Let $R$ be a discrete valuation
ring over $\CC$, let $K$ be its fraction field, and let
$(C_K,p_1,\ldots,p_n)$ be a stable pointed curve over $K$.
Think of $\operatorname{Spec}K$ as the punctured base and of
the closed point of $\operatorname{Spec}R$ as the missing limit.
The \emph{pointed stable reduction theorem} supplies, after a
finite extension of the fraction field and a corresponding
extension of discrete valuation rings $R\subset R'$, a stable
pointed family $C'\to\operatorname{Spec}R'$ whose generic fiber
is the prescribed base change of $C_K$.

The second external input is \emph{uniqueness of stable limits}:
if $C/R$ and $D/R$ are two stable pointed families, then every
specified pointed isomorphism $C_K\to D_K$ extends to a unique
pointed isomorphism $C\to D$. The word ``specified'' matters.
An object may have automorphisms; uniqueness concerns extension
of one given generic isomorphism, not uniqueness of all possible
isomorphisms between two limits.

Apply uniqueness to the valuative criterion for the representable
diagonal. The diagonal is already affine and of finite type, hence
quasi-compact and separated. The criterion proves that it is
proper, so $\Mgn$ is separated. In particular its diagonal is
finite, since an affine proper morphism is finite.
Apply existence to the valuative criterion for the finite type
stack $\Mgn\to\operatorname{Spec}\CC$. The permitted finite
base extension in this criterion is exactly what stable reduction
provides. The criterion proves universal closedness and hence
properness. The hypotheses of the criterion have all been supplied:
algebraicity, finite type, and quasi-separatedness came from the
Hilbert presentation, and separatedness was just checked.
Local noetherianness follows from finite type over $\CC$.

Why does stable reduction involve geometry beyond descent?
An \'etale cover sees every point of the base. The inclusion
$\operatorname{Spec}K\hookrightarrow\operatorname{Spec}R$
misses the closed point, so the stack condition cannot manufacture
a curve over that point. Stable reduction is the theorem that
does so, using the particular geometry of nodal stable curves.

\begin{Ex}[A genus-two stable boundary point]
Let $E_1,E_2$ be smooth elliptic curves and identify one point
on each to form $C=E_1\cup_qE_2$. Its arithmetic genus is two,
and each elliptic component has one node as a special point,
so $C$ is stable. A smoothing of this node has local equation
$xy=t$; the smooth fibers have genus two. The stable fiber at
$t=0$ belongs to $\overline{\mathcal M}_2$, while it does not
belong to the open smooth-curve stack $\mathcal M_2$.
The deformation theorem below supplies such smoothings. Uniqueness
of stable limits shows that this punctured family cannot instead
acquire a smooth special fiber after finite base extension.
This illustrates why $\mathcal M_2$ is not proper and why the
bar in $\overline{\mathcal M}_2$ matters.
\end{Ex}

\subsection{Smoothness and actual dimension}

\begin{Prop}\label{prop:curves-smooth}
The stack $\Mgn$ is smooth over $\CC$ and has pure dimension
$3g-3+n$.
\end{Prop}

\noindent\textbf{Original source and proof strategy.}
Deligne--Mumford \cite[Lemmas (1.3)--(1.4), Proposition (1.5),
pp.~79--82, and Theorem (5.2), pp.~104--105]{DM} analyze
deformations of nodal curves and obtain the unpointed smoothness
statement. Knudsen \cite[Theorem 2.7, p.~179]{Knudsen} gives the
pointed theorem. Alper \cite[Propositions 6.3.26 and 6.3.28,
pp.~232--235; Theorem 6.4.16, p.~241]{Alper} computes the
pointed deformation spaces and their dimension in detail.
We explain the count on the normalization and state the deformation
theorem used to pass from that count to smoothness.

\medskip
\noindent\textbf{Our guided proof.}
Smoothness of a Deligne--Mumford stack over $\CC$ can be tested
on an \'etale atlas. Its dimension is the dimension of that atlas
locally, because an \'etale map has relative dimension zero.
To prove smoothness we use the infinitesimal lifting criterion:
a small deformation over an Artinian base must lift when the base
is enlarged by a square-zero ideal. A deformation of a pointed
curve means a family together with an identification of its special
fiber with that fixed pointed curve.

Our \emph{external nodal deformation theorem} has two parts.
First, deformations of a proper nodal curve with distinct smooth
marked points have no obstruction to such infinitesimal lifting.
Second, the vector space of first-order deformations fits in
\begin{equation}\label{eq:curve-deformations}
 0\longrightarrow
 \bigoplus_v H^1(C_v,T_{C_v}(-B_v))
 \longrightarrow \operatorname{Def}(C,p_1,\ldots,p_n)
 \longrightarrow \bigoplus_{q\text{ a node}} T_q
 \longrightarrow0,
\end{equation}
where each $T_q$ is a one-dimensional smoothing space.
Choosing local coordinates at the two branches identifies it
with $K$, with deformation equation $xy=\epsilon a$.
That identification depends on coordinates, but its dimension
does not. The theorem is Alper's
\cite[Propositions 6.3.26 and 6.3.28, pp.~232--235]{Alper}.
The local calculation behind unobstructedness is that a node is
a single equation $xy=0$ whose deformation parameter can be
lifted, while a curve has no degree-two coherent cohomology to
obstruct the global gluing. The theorem makes that local-to-global
argument precise; we are not developing a general obstruction
theory in these notes.

The left side of \eqref{eq:curve-deformations} varies the smooth
normalization components, with every marked point and every
branch of a node remembered. Such variations keep the nodes
unsmoothed. The right side allows one smoothing parameter for
each node. Surjectivity says that these local smoothing parameters
can be realized by a global first-order deformation.

Let $V$ be the number of normalization components and let
$\delta$ be the number of nodes. Riemann--Roch and the vanishing
of $H^0(C_v,T_{C_v}(-B_v))$ proved above give
\[
 h^1(C_v,T_{C_v}(-B_v))=3h_v-3+m_v.
\]
There are $n$ marked points and two branches per node, so
$\sum_vm_v=n+2\delta$. The normalization sequence for functions
gives the genus formula
\[
                  g=\sum_vh_v+\delta-V+1.
\]
Taking dimensions in \eqref{eq:curve-deformations} now gives
\begin{align*}
 \dim\operatorname{Def}(C,p_1,\ldots,p_n)
 &=\sum_v(3h_v-3+m_v)+\delta\\
 &=3\sum_vh_v-3V+n+3\delta\\
 &=3g-3+n.
\end{align*}
There are no infinitesimal automorphisms, by
Proposition~\ref{prop:curves-dm}, so this is the tangent dimension
on an \'etale chart, without any extra infinitesimal symmetry
directions to remove. Unobstructedness and the infinitesimal
lifting criterion show that the charts are smooth; the constant
tangent dimension then proves the dimension assertion.

\begin{Ex}[Checking the dimension at a nodal genus-two curve]
For $C=E_1\cup_qE_2$, each normalization component is elliptic
with one special point. It contributes
$3\cdot1-3+1=1$ parameter. The node contributes one more.
The total is $1+1+1=3$, just as at a smooth genus-two curve.
The locus where this particular node remains unsmoothed has
only the first two parameters. A singular member of the family
does not force the moduli stack itself to be singular.
\end{Ex}

\begin{Ej}\label{ex:curve-dimension}
Glue two copies of $\mathbb P^1$ along three distinct pairs of
points, with no markings. Check stability and genus. Use
\eqref{eq:curve-deformations} to count the dimension of its
deformation space. Which parameters come from the components,
and which come from the nodes?
\end{Ej}

\subsection{What depends on the hypotheses}

All four propositions hold for every stable pair $g,n$ in the
setup of this chapter. The historical theorems actually prove
the corresponding proper smooth Deligne--Mumford statements over
$\operatorname{Spec}\mathbb Z$; we used their base change to
$\CC$, with the characteristic-zero expositions indicated.

If $2g-2+n\leq0$, the stable pointed moduli problem as defined
here has no geometric objects. One can still speak of its empty
stack, but not of a nonempty smooth stack of dimension $3g-3+n$.
If one instead drops stability while keeping those curves, the
conclusions change: an unmarked $\mathbb P^1$ has the
positive-dimensional group $\PGL_2$ of automorphisms, and an
unmarked elliptic curve has translations. Thus the argument for
the Deligne--Mumford property would fail.

Finally, smoothness belongs to the \emph{stack}; it is not an
automatic assertion about its coarse moduli space. Finite group
quotients of smooth charts can be singular. We have not needed
the existence or projectivity of a coarse moduli space for either
the stack proof or the properties proved here.

